% Ad Atticum 2.25 (ad-atticum-02-25) --- English translation
% Translated by Claude (Anthropic) from the Latin (Perseus).
% Style guide: STYLE.md.

\ciceroLetterOpener{CICERO ATTICO salutem}

\ciceroSection{1} When I have praised someone in your house, I shall want him to learn it from you that I have done so --- as lately you know I wrote to you of Varro's good service to me, and you wrote back that this gave you the highest pleasure. But I should have preferred that you had written to him that he was satisfying me --- not because he was doing so, but to make him do so. For he is wonderfully odd in his ways, as you know: \textit{[Greek: tortuous and never sound at all]}. But I hold to that maxim, \textit{[Greek: bear with the ways of the powerful]}. By Hercules, your other friend Hortalus --- with what full hand, how openly, how splendidly he lifted our praises to the stars when he spoke of Flaccus's praetorship and of that time of the Allobroges! Take this for fact: nothing could have been said more affectionately, more honourably, more abundantly. I should much like you to write him that I have sent this report of him to you.

\ciceroSection{2} But why should you write? You whom I now think to be on the way and present; for so I dealt with you in my last letter. I look for you eagerly, miss you eagerly, and I no more so than the very matter and the moment demand. What shall I write to you about these affairs except what I have often written? Of the commonwealth nothing more desperate; of those by whose work this is, nothing in greater hatred. We are fortified, as our opinion and hope and conjecture have it, by the firmest goodwill of men. Therefore, fly here. Either you will free us of every trouble or you will be partaker. I am the briefer for this reason: that, as I hope, in a short time we shall be able to confer face to face about what we wish. Take care of your health.
